\section*{Question 4}

\subsection*{(a) What is the major mutation signature contribute to AML?}

According to the given matrix, AML is associated with \textbf{Signature 1} and \textbf{Signature 5}.

\boxed{\textbf{Signature 1}} is the most dominant across all cancer types and typically shows stronger contribution in AML.

\subsection*{(b) What is the major mutation type related to AML?}

\boxed{\textbf{C>T}} is accounted for more than 10\% of all.

\subsection*{(c) Signatures Dominated by T>A}

\boxed{\textbf{Signature 22 and Signature 27}} are dominated by T>A, which is more than 10\%.

\textbf{Associated Cancer Types}: Liver cancer, Urothelial carcinoma (Signature 22); Kidney clear cell carcinoma (Signature 27).

\subsection*{(d) Model Design}

\subsubsection*{Feature Representation}
Each patient can be represented by:
\begin{itemize}
    \item \textbf{Raw mutation spectrum} (96-dimensional trinucleotide mutation frequencies), or
    \item \textbf{Signature contributions} obtained from NMF ($H$ vector).
\end{itemize}

\subsubsection*{Model Choice}
As far as the task is concerned, we can use a two-step model:

\[
    \textbf{NMF decomposition} \; \Rightarrow \; \textbf{Random Forest classifier}
\]

\subsubsection*{Procedure}
\begin{enumerate}
    \item Build matrix $X$ from patient mutation profiles.
    \item Apply NMF to obtain $H$ (signature weights per patient).
    \item Split the data into training and test sets.
    \item Train a Random Forest or Logistic Regression model using $H$ as input and cancer type as target.
    \item For a new patient, compute mutation profile $\rightarrow$ project to $H$ space $\rightarrow$ predict cancer type.
\end{enumerate}

\subsubsection*{Why This Model Fits the Task}
\begin{itemize}
    \item NMF provides biologically interpretable latent features (mutation signatures).
    \item The classifier captures statistical relationships between signature patterns and cancer types.
    \item The model can generalize to unseen patients with similar mutational patterns.
    \item Probability outputs allow risk estimation and uncertainty quantification.
\end{itemize}